\chapter{Permission before execution}\label{ch:permission}
The schema in the preceding chapter tells us what the world contains.
It does not tell us which proposed movement the world permits. This missing
step is easy to overlook when a quoted action has an attractive value.
An evaluator can calculate a field difference for a negative-stock endpoint
just as readily as for a valid one. Software must therefore distinguish
the ability to calculate a number from authority to perform its movement.

\section{Three questions, three records}
First ask whether the declared physical mechanism can perform the action.
Second ask what the action would do to the represented potential. Third,
if the prospective account policy has been adopted, ask whether each owner
can afford its assigned change. An action can pass one question and fail
another. Keeping separate records lets a reviewer understand why it was
refused without pretending that every refusal expresses one universal
notion of safety.

Physical permission includes the source and destination constraints,
allowed topology, carrier identity and timing. Field valuation depends on
the frozen state and potential. Account permission depends additionally on
the owner map and prior balances. The physical-random comparator omits the
third filter; it does not get a different tank capacity or a privileged
right to create resource.

\section{Endpoint checks are necessary, not a complete plant}
For a closed lossless action, a simple check verifies
\begin{equation}
 z+\delta_G\geq0,\qquad \mathbf1^T\delta_G=0.
\end{equation}
It is tempting to call this the whole permission rule. Yet a pipe can have
a finite delivery rate, a machine can serve only one request at a time, and
a destination can have a storage ceiling. The first domain may intentionally
omit these complications. It must say so. An omitted constraint is not a
constraint that the software has proved satisfied.

For a declared constant-rate aggregate path, convexity of the nonnegative
simplex keeps intermediate stocks nonnegative when both endpoints are valid.
That argument concerns aggregate coordinates. It does not authorize an
individual machine's rate or prove the causal availability of incoming
material to an outgoing child action. The path and plant declarations must
agree about what simultaneity means.

\section{A supplied budget is not a permission theorem}
Suppose a mechanism is given two requests of four units and a joint budget
of six. Proportional scaling gives three units to each. This arithmetic can
be reused wherever proportional scaling is the declared resolver. It says
nothing about why the budget is six. A source-preservation certificate,
a physical throughput limit and a deliberate experimental restriction can
all supply the same numerical budget while making different claims.

The distinction is especially important for historical P1C. Its robust
budget routine assumes an eligible regenerative source. Applying the
formula to a finite store without that classification does not inherit the
certificate. Equally, applying the old classifier's finite-extraction
prohibition to every internal relocation would prevent the new closed world
from expressing its intended question. Neither change should be smuggled
through a helper function with a familiar name.

\section{Export protection is not destination recovery}
The historical quantity $R_{\mathrm{eff}}$ belongs to a source export
certificate. A destination's old homeostatic reserve $R$ belongs to the
potential and condition of that destination. Even when a source budget
protects its own boundary perfectly, it does not choose where to send stock
or guarantee that another node's reserve is restored.

In the active Gaussian domain, the reference $x^*$ is neither of these
objects. It defines a comparison geometry, not a prohibition against
physically crossing a reference value. Otherwise the physical permission
layer would become a reference-seeking controller before the actor made a
choice, and the proposed comparison would test a different mechanism.

\section{Resolve first, then bind the exact vector}
If a declared resolver modifies proposals, record both the requests and
the accepted quantities. The evaluator receives the accepted vector and
its complete support. The eventual execution must apply that same vector,
not a newly clipped version with the old value attached.

If a late physical condition invalidates execution, a prospective system
needs an explicit refusal or re-epoching procedure. It may not retain the
old positive attribution while silently reducing deliveries. Whether to
retry is also a policy decision: repeated attempts can change selection
probabilities and consume different random draws.

\section{A permission review in plain language}
For each refusal, a reader should be able to finish one sentence:
``This proposal could not proceed because this named requirement failed
against this version of the state.'' A missing owner, a negative stock and
an unavailable execution permit are different failures. Mixing them into
one generic false value makes debugging easier only by making the science
harder to audit.

The historical source reading that follows shows why the premises of an
actual permission mechanism must remain beside its formula. After it,
the event-order chapter will place these records into one explicit chronology.
